%-------------------------
% Resume in LaTeX - Exact match to RANGARAJAN_RESUME.docx
% Page: A4, Margins: top=0.138in, right/bottom/left=0.099in
% Font: Times New Roman 10pt
%-------------------------

\documentclass[a4paper,10pt]{article}

\usepackage[T1]{fontenc}
\usepackage{mathptmx} % Times New Roman
\usepackage[empty]{fullpage}
\usepackage{enumitem}
\usepackage[hidelinks]{hyperref}
\usepackage{tabularx}
\usepackage{geometry}
\usepackage{titlesec}
\usepackage{fancyhdr}

\geometry{
  a4paper,
  top=0.138in,
  bottom=0.099in,
  left=0.099in,
  right=0.099in
}

\pagestyle{empty}
\setlength{\parindent}{0pt}
\setlength{\parskip}{0pt}
\raggedbottom
\raggedright

% Section heading: bold, uppercase, with full-width rule below
\titleformat{\section}{
  \vspace{2pt}\bfseries\raggedright\normalsize
}{}{0em}{}[\vspace{-6pt}\rule{\linewidth}{0.4pt}\vspace{-4pt}]

\titlespacing{\section}{0pt}{6pt}{4pt}

% Bullet list settings to match DOCX
\setlist[itemize]{
  leftmargin=0.18in,
  itemsep=0pt,
  parsep=0pt,
  topsep=2pt,
  partopsep=0pt
}
\renewcommand{\labelitemi}{$\bullet$}

%-------------------------------------------
\begin{document}

%---------- HEADING ----------
\begin{center}
  {\large \textbf{RANGARAJAN NADADHUR RAMESH}} \\[2pt]
  {\small rangarajannr27@gmail.com \textbar{} +16027488247 \textbar{} https://www.linkedin.com/in/rangarajan-n-r}
\end{center}

\vspace{2pt}

%---------- SUMMARY ----------
\section{SUMMARY}

Accomplished Design Engineer with expertise in SolidWorks, ANSYS FEA, GD\&T, and DFM, driving mechanical product development from concept to production.
Proven ability to optimize designs, reduce costs, and improve manufacturing efficiency through analytical problem-solving and cross-functional collaboration.

%---------- EDUCATION ----------
\section{EDUCATION}

\noindent \textbf{Master of Science in Mechanical Engineering} \hfill \textit{May 2026} \\
Arizona State University, Arizona, USA \hfill \textit{GPA: 3.70}

\vspace{2pt}
\noindent \textbf{Relevant Coursework:} Modern Manufacturing Methods, Intermediate Thermodynamics

\vspace{4pt}

\noindent \textbf{Bachelor of Technology in Mechanical Engineering} \hfill \textit{May 2021} \\
Dayananda Sagar University, Karnataka, India \hfill \textit{CGPA: 8.36}

\vspace{2pt}
\noindent \textbf{Additional Coursework:} GD\&T based on ASME Y14.5 (Udemy), Lean Six Sigma Yellow Belt (Udemy)

%---------- TECHNICAL SKILLS ----------
\section{TECHNICAL SKILLS}

\begin{itemize}[leftmargin=0.15in, itemsep=0pt, parsep=0pt, topsep=2pt]
  \item \textbf{CAD Modeling:} SolidWorks, Siemens NX, PTC Creo, Fusion 360, AutoCAD
\item \textbf{Simulation:} ANSYS FEA, ANSYS Fluent, Tolerance Stack-Up, GD&T (ASME Y14.5), FMEA
\item \textbf{Engineering & Documentation:} DFM, DFA, BOM Management, Engineering Change Documentation, NPD
\item \textbf{Manufacturing:} Plastic and Sheet Metal Design, Twin-Sheet Thermoforming, 3D Printing, CNC
\end{itemize}

%---------- PROFESSIONAL EXPERIENCE ----------
\section{PROFESSIONAL EXPERIENCE}

\noindent \textbf{Aktis Engineering Solutions Pvt. Ltd., Bangalore, India \textbar{} Design Engineer} \hfill \textit{Sep 2022 to Jul 2024}
\begin{itemize}
  \item Designed production-ready plastic and sheet metal components and assemblies in SolidWorks, generating 3D CAD models and GD&T drawings.
\item Improved production release quality by 55% through structured validation workflows and PFMEA contributions.
\item Reduced fit and assembly issues by 30% by integrating supplier feedback into tolerance stack-up analysis during DFM/DFA reviews.
\item Achieved 15% material cost reduction and 30% to 50% performance improvement on twin-sheet thermoformed components using ANSYS FEA.
\end{itemize}

\vspace{2pt}
\noindent \textbf{Pactera EDGE (now Centific), Hyderabad, India \textbar{} Project Associate} \hfill \textit{Mar 2022 to Aug 2022}
\begin{itemize}
  \item Improved team output quality from 75% to 99% and individual accuracy to 95% by identifying workflow gaps and implementing process controls.
\end{itemize}

\vspace{2pt}
\noindent \textbf{Bosch Rexroth in Collaboration with DSU, Bangalore, Karnataka, India \textbar{} Intern} \hfill \textit{Feb 2021 to Apr 2021}
\begin{itemize}
  \item Developed hydraulic and pneumatic circuit layouts and PLC ladder logic for prototype industrial automation equipment.
\end{itemize}

%---------- ACADEMIC PROJECTS ----------
\section{ACADEMIC PROJECTS}

\noindent \textbf{Study of Twin-Sheet Thermoforming}
\begin{itemize}[topsep=1pt]
  \item Designed Twin-Sheet structures in SolidWorks and performed FEA using ANSYS, analyzing deformation and stiffness.
  \item Demonstrated structural and economic advantages of Twin-Sheet Thermoforming compared to conventional processes.
  \item Led a 4-member team studying Twin-Sheet Thermoforming and its advantages over traditional manufacturing processes.
\end{itemize}
\vspace{2pt}
\noindent \textbf{Honeycomb Core Sandwich Composite Panels - Bending Properties}
\begin{itemize}[topsep=1pt]
  \item Designed and fabricated lightweight honeycomb core sandwich composite panels using Fusion 360 for bending strength analysis.
  \item Conducted 2-way ANOVA (DoE) to optimize core geometry configurations, enhancing structural properties.
  \item Managed 3D printing and material selection to evaluate prototype performance in terms of strength-to-weight.
\end{itemize}
\vspace{2pt}
\noindent \textbf{Self-Propelled Onion Harvester}
\begin{itemize}[topsep=1pt]
  \item Led a 25-member team developing an onion harvester, reducing maintenance cost by 40% and harvesting time by 50%.
  \item Contributed to chassis design, steering mechanism, and powertrain gear ratio development for improved efficiency.
  \item Engineered a laser and image-detection-based yield estimation solution, reducing estimation time by 50%.
\end{itemize}
\vspace{2pt}
\noindent \textbf{Optimizing Wind Turbine Operation for Bird Collision Prevention}
\begin{itemize}[topsep=1pt]
  \item Led a 3-member team developing a wind turbine collision avoidance system to reduce bird strikes.
  \item Designed a multi-sensor detection system integrating radar and LiDAR for real-time object detection.
  \item Showcased multi-million-dollar cost savings potential by reducing bird collisions and associated fines.
\end{itemize}
\vspace{2pt}
\noindent \textbf{Study of Two-Bounce LiDAR for Non-Line-of-Sight Imaging}
\begin{itemize}[topsep=1pt]
  \item Collaborated in a 5-member team studying Two-Bounce LiDAR and NLOS imaging for autonomous vehicle applications.
  \item Assessed non-line-of-sight object detection capability of SPAD sensors, lasers, and optics.
  \item Presented feasibility, risks, deployment recommendations, and economic considerations for real-world implementation.
\end{itemize}

\end{document}
